% FAC-EN-U0078
% Source: fac_body.tex lines 3295--3356.
% Authority: PDF physical pp. 80--81; printed pp. 275--276
% (no. 80 complete; hard stop before no. 81).

\medskip
\phantomsection\label{fac:III:6:no:80}
\textbf{80. Relation with the characteristic function of a graded
$S$-module.}
Let $\mathfrak F$ be a coherent algebraic sheaf on the space
$\mathbf P_r(K)$; we shall write $\chi(\mathfrak F)$ instead of
$\chi(\mathbf P_r(K),\mathfrak F)$. We have:

\phantomsection\label{fac:III:6:no:80:prop:3}
\textsc{Proposition 3.} \textit{$\chi(\mathfrak F(n))$ is a polynomial
in $n$ of degree at most $r$.}

By \hyperref[fac:III:2:no:60:theorem:2]{Theorem~2 of no.~60}, there
exists a finitely generated graded $S$-module $M$ such that
$\mathfrak A(M)$ is isomorphic to $\mathfrak F$. Applying Hilbert's
syzygy theorem to $M$, we obtain an exact sequence of graded
$S$-modules
\phantomsection\label{fac:III:6:no:80:prop:3:module-resolution}
\[
0\longrightarrow L^{r+1}\longrightarrow\cdots\longrightarrow L^0
\longrightarrow M\longrightarrow0,
\]
where the $L^q$ are finitely generated free modules. Applying the
functor $\mathfrak A$ to this sequence, we obtain\footnote{The French
print omits the subject \emph{on} in ``En appliquant le foncteur
$\mathfrak A$ \`a cette suite, obtient une suite exacte de faisceaux.''
The diplomatic French layer preserves the printed sentence; the
corrected French layer restores \emph{on}, and the English follows the
grammatically and mathematically required reading.} an exact sequence
of sheaves
\phantomsection\label{fac:III:6:no:80:prop:3:sheaf-resolution}
\[
0\longrightarrow\mathcal L^{r+1}\longrightarrow\cdots
\longrightarrow\mathcal L^0\longrightarrow\mathfrak F
\longrightarrow0,
\]
where each $\mathcal L^q$ is isomorphic to a finite direct sum of
sheaves $\mathcal O(n_i)$. \hyperref[fac:III:6:no:79:prop:2]{Proposition~2}
shows that $\chi(\mathfrak F(n))$ is equal to the alternating sum of
the $\chi(\mathcal L^q(n))$, which reduces us to the case of the sheaf
$\mathcal O(n_i)$. But it follows from
\hyperref[fac:III:3:no:62]{no.~62} that
\phantomsection\label{fac:III:6:no:80:formula:chi-O-n}
\[
\chi(\mathcal O(n))=\binom{n+r}{r},
\]
which is indeed a polynomial in $n$ of degree at most $r$; this proves
the \hyperref[fac:III:6:no:80:prop:3]{proposition}.

\phantomsection\label{fac:III:6:no:80:prop:4}
\textsc{Proposition 4.} \textit{Let $M$ be a graded $S$-module
satisfying condition \textup{(TF)}, and let
$\mathfrak F=\mathfrak A(M)$. For every sufficiently large $n$,}
\phantomsection\label{fac:III:6:no:80:prop:4:formula}
\[
\chi(\mathfrak F(n))=\dim_K M_n.
\]

Indeed, we know from \hyperref[fac:III:3:no:65]{no.~65} that, for
sufficiently large $n$, the homomorphism
$\alpha:M_n\to H^0(X,\mathfrak F(n))$ is bijective and
$H^q(X,\mathfrak F(n))=0$ for every $q>0$; hence
\[
\chi(\mathfrak F(n))=h^0(X,\mathfrak F(n))=\dim_K M_n.
\]

We thus recover the well-known fact that $\dim_K M_n$ is a polynomial
in $n$ for sufficiently large $n$. This polynomial, which we shall
denote by $P_M$, is called the
\phantomsection\label{fac:III:6:no:80:def:characteristic-function}
\textit{characteristic function} of $M$. For every $n\in\mathbf Z$,
we have
\phantomsection\label{fac:III:6:no:80:formula:P-M}
\[
P_M(n)=\chi(\mathfrak F(n));
\]
in particular, on taking $n=0$, we see that
\textit{the constant term of $P_M$ is equal to $\chi(\mathfrak F)$}.

Apply \hyperref[fac:III:6:no:80:prop:4]{this} to $M=S/I(V)$, where $I(V)$ is the homogeneous ideal of $S$
consisting of the polynomials that vanish on a closed subvariety $V$
of $\mathbf P_r(K)$. In this case the constant term of $P_M$ is called
the \phantomsection\label{fac:III:6:no:80:def:arithmetic-genus}
\textit{arithmetic genus} of $V$ (cf.~[19]); since, on the other hand,
$\mathfrak A(M)=\mathcal O_V$, we obtain:

\phantomsection\label{fac:III:6:no:80:prop:5}
\textsc{Proposition 5.} \textit{The arithmetic genus of a projective
variety $V$ is equal to}
\phantomsection\label{fac:III:6:no:80:prop:5:formula}
\[
\chi(V,\mathcal O_V)=\sum_{q=0}^{\infty}(-1)^q
 \dim_K H^q(V,\mathcal O_V).
\]

\textsc{Remarks}. \phantomsection\label{fac:III:6:no:80:remark:1}
(1) The \hyperref[fac:III:6:no:80:prop:5]{preceding proposition} brings out the fact that the arithmetic
genus is independent of the embedding of $V$ in a projective space,
since the same is true of the $H^q(V,\mathcal O_V)$.

\phantomsection\label{fac:III:6:no:80:remark:2}
(2) The \textit{virtual} arithmetic genus (defined by Zariski in [19])
can likewise be reduced to an Euler--Poincaré characteristic. We shall
return later to this question, which is closely connected with the
Riemann--Roch theorem.

\phantomsection\label{fac:III:6:no:80:remark:3}
(3) For convenience, we have adopted a definition of the arithmetic
genus slightly different from the classical definition (cf.~[19]). If
all the irreducible components of $V$ have the same dimension $p$, the
two definitions are related by the formula
\phantomsection\label{fac:III:6:no:80:remark:3:formula}
\[
\chi(V,\mathcal O_V)=1+(-1)^p p_a(V).
\]
